\begin{lemma}
\label{lemma-ffdescent-ML}
\begin{reference}
Email from Juan Pablo Acosta Lopez dated 12/20/14.
\end{reference}
Let $R \to S$ be a faithfully flat ring map. Let $M$ be an $R$-module. If the
$S$-module $M \otimes_R S$ is Mittag-Leffler, then $M$ is Mittag-Leffler.
\end{lemma}

\begin{proof}
Write $M = \colim_{i\in I} M_i$ as a directed colimit
of finitely presented $R$-modules $M_i$.
Using Proposition \ref{proposition-ML-characterization}, we see that we
have to prove that for each $i \in I$ there exists $i \leq j$, $j\in I$
such that $M_i\rightarrow M_j$ dominates $M_i\rightarrow M$.

\medskip\noindent
Take $N$ the pushout
$$
\xymatrix{
M_i  \ar[r] \ar[d] & M_j \ar[d] \\
M \ar[r] & N
}
$$
Then the lemma is equivalent to the existence of $j$ such that
$M_j\rightarrow N$ is universally injective, see
Lemma \ref{lemma-domination-universally-injective}.
Observe that the tensorization by $S$
$$
\xymatrix{
M_i\otimes_R S  \ar[r] \ar[d] & M_j\otimes_R S \ar[d] \\
M\otimes_R S \ar[r] & N\otimes_R S
}
$$
Is a pushout diagram.  So because
$M \otimes_R S = \colim_{i\in I} M_i \otimes_R S$
expresses $M\otimes_R S$ as a colimit of $S$-modules of
finite presentation, and $M\otimes_R S$ is Mittag-Leffler,
there exists $j \geq i$ such that $M_j\otimes_R S\rightarrow N\otimes_R S$
is universally injective. So using that $R\rightarrow S$ is faithfully flat
we conclude that $M_j\rightarrow N$ is universally injective too.
\end{proof}
